% ============================================================
% Come si usano, argomento 7: varianza e disuguaglianze
% ============================================================
\argomento{Varianza e disuguaglianze}{Formule a p.~\pageref*{f:rms}}

\begin{scelta}{Come scegliere la formula}
\renewcommand{\arraystretch}{1.25}
\begin{tabular}{@{}P{0.58\linewidth}P{0.36\linewidth}@{}}
  \toprule
  \textbf{Se il testo\ldots} & \textbf{si usa\ldots} \\
  \midrule
  chiede quanto la variabile è dispersa, la varianza, la deviazione standard
    & \vv{varianza}, quasi sempre con $E[X^2] - \mu_X^2$ \vv{varrms} \\
  parla di potenza media o di valore efficace & il valore quadratico medio \vv{rms} \\
  sposta o riscala la variabile (unità di misura, guadagno, offset) & \vv{lineare} \\
  dà solo la media di una variabile non negativa e chiede un limite a $P(Z \geq \delta)$ & Markov \vv{markov} \\
  dà media e deviazione standard e chiede un limite alla probabilità di allontanarsi dalla media, o di
    restarle vicino & Chebyshev \vv{chebyshev} \\
  \bottomrule
\end{tabular}

Le varianze delle variabili notevoli sono nelle loro voci: Bernoulli \vv{bernoulli}, binomiale \vv{distbinomiale},
uniforme \vv{uniforme}, Poisson \vv{poisson}.
\end{scelta}

\scheda{rms}{
  \usoquando{Si vuole la potenza media di un segnale (la potenza dipende dal quadrato) o il valore efficace di una
    tensione; oppure come passo intermedio per la varianza (\vv{varrms}).}
  \usocome{Si pesa il quadrato di ogni valore con la sua probabilità: è il teorema fondamentale \vv{valoremedio}
    con $g(x) = x^2$.}
  \usoesempio{$X$ vale 4, 5, 6 con probabilità $\frac{1}{4}, \frac{1}{2}, \frac{1}{4}$:
    $E[X^2] = 16 \cdot \frac{1}{4} + 25 \cdot \frac{1}{2} + 36 \cdot \frac{1}{4} = 25{,}5$ e
    $X_{rms} = \sqrt{25{,}5} \simeq 5{,}05$.}
}

\scheda{varianza}{
  \usoquando{Si vuole misurare quanto la variabile si allontana dalla sua media. La deviazione standard ha la stessa
    unità di misura di $X$, ed è quella da leggere.}
  \usocome{Si media il quadrato dello scarto dalla media. Di solito è più rapido passare da $E[X^2]$ \vv{varrms}.}
  \usoesempio{Due variabili con media 5: $X \in \{4, 5, 6\}$ con probabilità $\frac{1}{4}, \frac{1}{2},
    \frac{1}{4}$ ha $\sigma_X^2 = \frac{1}{2}$ ($\sigma_X \simeq 0{,}71$); $Y \in \{1, 5, 9\}$ con le stesse
    probabilità ha $\sigma_Y^2 = 8$ ($\sigma_Y \simeq 2{,}83$).}
}

\scheda{varrms}{
  \usoquando{Quasi sempre, per calcolare una varianza: $E[X^2]$ e $\mu_X$ si trovano separatamente.}
  \usoattenzione{$E^2[X]$ vuol dire $(E[X])^2$, non $E[X^2]$. Una varianza negativa segnala un errore nei conti.}
  \usoesempio{Per la $X$ dell'esempio precedente: $\sigma_X^2 = 25{,}5 - 5^2 = 0{,}5$.}
}

\scheda{momenti}{
  \usoquando{Serve $E[X^2]$ di una binomiale o di una Poisson, per esempio per ricavarne la varianza.}
  \usocome{Conviene calcolare prima $E[X(X-1)]$, perché il fattore $k(k-1)$ si semplifica con $k!$; poi si aggiunge
    $E[X]$.}
  \usoesempio{$X \sim \mathcal{B}(10;\ 0{,}1)$: $E[X(X-1)] = 10 \cdot 9 \cdot 0{,}01 = 0{,}9$, $E[X^2] = 0{,}9 + 1 =
    1{,}9$ e $\sigma_X^2 = 1{,}9 - 1^2 = 0{,}9 = N p\,(1-p)$.}
}

\scheda{lineare}{
  \usoquando{La variabile viene spostata o riscalata: cambi di unità di misura, guadagni e offset, standardizzazione.}
  \usocome{Una costante sommata sposta la media e non cambia la varianza; un fattore moltiplica la media e, al
    quadrato, la varianza.}
  \usoattenzione{$\sigma_{aX+b}^2 = a^2 \sigma_X^2$, non $a\, \sigma_X^2$; con $a = -1$ la varianza non cambia.}
  \usoesempio{Una temperatura in gradi Celsius ha $\mu = 20$ e $\sigma = 2$. In Fahrenheit, $F = 1{,}8\, C + 32$:
    $\mu_F = 1{,}8 \cdot 20 + 32 = 68$ e $\sigma_F = 1{,}8 \cdot 2 = 3{,}6$. Un dado onesto è l'uniforme su
    $\{0, \ldots, 5\}$ più 1, quindi ha la sua stessa varianza, $\frac{35}{12}$.}
}

\scheda{segno}{
  \usoquando{Come controllo dei risultati: una variabile che non è mai negativa non può avere media negativa, e
    nessuna varianza può essere negativa.}
  \usoesempio{Da $\sigma_X^2 \geq 0$ segue $E[X^2] \geq E^2[X]$: una variabile con $E[X] = 3$ ed $E[X^2] = 8$ non
    può esistere.}
}

\scheda{markov}{
  \usoquando{Di una variabile non negativa si conosce solo la media, e si vuole un tetto alla probabilità che assuma
    valori grandi.}
  \usoattenzione{Vale solo per variabili non negative, e dà un limite superiore spesso molto largo.}
  \usoesempio{I pacchetti in coda in un router sono in media 2: $P(Z \geq 10) \leq \frac{2}{10} = 0{,}2$. Se si
    sapesse che $Z$ è di Poisson, il valore esatto sarebbe circa $0{,}00005$.}
}

\scheda{chebyshev}{
  \usoquando{Si conoscono media e deviazione standard, ma non la pmf, e si vuole un limite alla probabilità di
    allontanarsi dalla media, o di restarle vicino.}
  \usocome{Si esprime la distanza dalla media in deviazioni standard: $k$ = distanza$/\sigma_X$.}
  \usovarianti{Con la distanza $\varepsilon$ scritta direttamente il limite è $\sigma_X^2 / \varepsilon^2$: è la forma
    che serve per la legge dei grandi numeri.}
  \usoattenzione{Per $k \leq 1$ non dice nulla, perché $1/k^2 \geq 1$. Vale per qualsiasi pmf, quindi spesso è largo.}
  \usoesempio{$\mu_X = 10$, $\sigma_X = 2$: $P(6 \leq X \leq 14) \geq 1 - \frac{1}{4} = 0{,}75$ (con $k = 2$) e
    $P(\abs{X - 10} \geq 6) \leq \frac{1}{9} \simeq 0{,}11$ (con $k = 3$).}
}

\scheda{concentrazione}{
  \usoquando{Per giudicare quanto una variabile è prevedibile rispetto alla sua grandezza: se $\mu_X / \sigma_X$ è
    grande, le fluttuazioni sono piccole in proporzione.}
  \usoesempio{Poisson con $\lambda = 100$: $\mu_X / \sigma_X = 10$, la variabile vale circa $100 \pm 10$
    (fluttuazioni del 10\%); con $\lambda = 10\,000$ vale circa $10\,000 \pm 100$ (l'1\%).}
}
